\documentclass{article}

\usepackage[utf8]{inputenc}
\usepackage[T1]{fontenc}
\usepackage{helvet}
\usepackage{geometry}
\usepackage{xcolor}
\usepackage{eso-pic}
\usepackage{fancyhdr}
\usepackage{parskip}
\usepackage{microtype}
\usepackage[english, ukrainian]{babel}
\usepackage{url}
\usepackage{amsmath}
\usepackage{amssymb}
\usepackage{amsthm}
\usepackage{longtable}
\usepackage{booktabs}
\usepackage{hyperref}
\usepackage{titlesec}
\usepackage{listings}

\definecolor{nato}{HTML}{293757}
\definecolor{footergray}{gray}{0.65}

\geometry{a4paper,left=3cm,right=3cm,top=3cm,bottom=3cm}
\renewcommand{\familydefault}{\sfdefault}
\setcounter{secnumdepth}{3}

\titleformat{\section}
  {\color{nato}\Huge\bfseries}{\thesection.}{1em}{}
\titlespacing*{\section}{0pt}{30pt}{12pt}

\titleformat{\subsection}
  {\color{nato}\LARGE\bfseries}{\thesubsection.}{1em}{}
\titlespacing*{\subsection}{0pt}{20pt}{8pt}

\titleformat{\subsubsection}
  {\color{nato}\Large\bfseries}{\thesubsubsection.}{1em}{}
\titlespacing*{\subsubsection}{0pt}{10pt}{6pt}

\fancyhf{}
\fancyhead[R]{\small\color{footergray} 18 квітня 2026}
\fancyfoot[L]{\small\color{footergray} Zen Crypted}
\fancyfoot[R]{\small\color{footergray}\thepage}

\newcommand\HUGE[1]{\fontsize{40}{40}\selectfont #1}

\renewcommand{\headrulewidth}{0pt}
\renewcommand{\footrulewidth}{0.4pt}
\renewcommand{\footrule}{\color{footergray}\hrule width \textwidth height 0.4pt}

\lstset{
  basicstyle=\ttfamily\small,
  breaklines=true,
  frame=single,
  backgroundcolor=\color{gray!5},
  columns=fullflexible
}

\begin{document}

\pagestyle{fancy}

\begin{titlepage}
  \AddToShipoutPictureBG*{\AtPageLowerLeft{\color{nato}\rule{\paperwidth}{\paperheight}}}
  \color{white}\thispagestyle{empty}\vspace*{4cm}\centering
  {\Huge  \textbf{MSC Z.120} \par} \vspace{2cm}
  {\HUGE  \textbf{Тестувальні сценарії \\ CHAT V.420.2} \par} \vspace{2.5cm}
  \vfill
  {\large 2026 \copyright\ Zen Crypted \par}
\end{titlepage}

\newpage

\tableofcontents

\vspace{3cm}

\section{Вступ}

Набір сценаріїв задає поведінку системи та використовується для перевірки її узгодженості. Кожен сценарій описує послідовність дій, початковий стан (given) та очікувані результати (expect), що дозволяє систематично перевіряти семантику протоколу Buddha на рівні kernel-моделі.

\section{Архітектура тестового набору}

Сценарії організовані як система тестів, що перевіряє семантику kernel-моделі та поведінку протоколу. Вони розбиті на тематичні групи відповідно до основних доменів: повідомлення, читання, відтворення, модерація, групи, пагінація, присутність, home, контакти, згадки, видимість та розширення. Кожен сценарій є самодостатнім і може виконуватися незалежно, забезпечуючи повне покриття invariants між protocol state, read/replay, moderation та ABAC visibility.

\section{Сценарії домену}

Сценарії домену описують основну поведінку системи та покривають ключові можливості протоколу. Вони перевіряють базові операції над повідомленнями, стрімами подій, курсорами та станом ресурсів.

\section{Складні сценарії та узгодженість}

Складні сценарії перевіряють взаємодію подій, порядок їх обробки та узгодженість системи в нетривіальних випадках. Вони включають reorder edit/delete, late delete, ban після прийнятого повідомлення, snapshot після видалення групи та federated convergence.

\section{Сценарії повідомлень}

Сценарії повідомлень перевіряють структуру даних, передачу повідомлень та операції над їх вмістом.

\begin{lstlisting}[language={},basicstyle=\ttfamily\small]
scenario structured send + structured expect
  session alice1
  connect
  auth

  session bob1
  connect
  auth

  session alice1
  send message to bob {
    body: "hi"
    subject: "Draft"
    priority: high
  }

  session bob1
  expect message from alice {
    body: "hi"
    subject: "Draft"
    priority: high
  }
\end{lstlisting}

(Повний набір сценаріїв з DSL-PAYLOAD.md: structured payload, validation rules, edit/delete convergence, given payload, replay convergence, mutation by captured/seeded id тощо.)

\section{Сценарії читання}

Сценарії читання перевіряють курсори, стан прочитаного та узгодженість поведінки read/unread.

\begin{lstlisting}[language={},basicstyle=\ttfamily\small]
scenario delivery + read

session alice
connect
auth

session bob
connect
auth

session alice
send message to bob "hi"

session bob
expect message from alice body "hi"

session bob
send read for last

session alice
expect message marked as read
\end{lstlisting}

(Повний набір з DSL-READ.md: basic delivery, read cursor, multi-session sync, read backward, read after reconnect, multi-feed isolation, partial read, unread boundary тощо.)

\section{Сценарії відтворення}

Сценарії відтворення перевіряють відновлення подій, межі відтворення, усунення пропусків та перехід між знімком і потоком подій.

\begin{lstlisting}[language={},basicstyle=\ttfamily\small]
scenario home bootstrap then replay

session alice
connect
auth

session bob
connect
auth
add alice to roster

session alice
send message to bob "m1"
send message to bob "m2"

session bob
disconnect
wait 500ms
reconnect

bootstrap home limit 20 preview 1

expect feeds
expect previews
expect shared snapshot

query events peer alice after snapshot

expect no duplicates
expect no gaps
\end{lstlisting}

(Повний набір з DSL-REPLAY.md: replay, preview, gap recovery, paged snapshot, multi-feed isolation, duplicate delivery тощо.)

\section{Сценарії модерації}

Сценарії модерації перевіряють обмеження дій, доступу та видимості без зміни канонічної історії подій.

\begin{lstlisting}[language={},basicstyle=\ttfamily\small]
scenario ban user

session alice
connect
auth

ban bob

expect bob is banned
\end{lstlisting}

(Повний набір з DSL-MODERATION.md: global/group ban, unban, moderation list, no rewrite history, block future actions тощо.)

\section{Сценарії груп}

Сценарії груп перевіряють групові стріми подій, склад учасників, ролі та правила доступу.

\begin{lstlisting}[language={},basicstyle=\ttfamily\small]
scenario create group

session alice
connect
auth

create group room1

expect group room1 exists
expect alice is owner of group room1
expect alice is member of group room1
\end{lstlisting}

(Повний набір з DSL-GROUP.md: create, add/remove member, send/query, delete group тощо.)

\section{Сценарії пагінації}

Сценарії пагінації перевіряють продовження вибірки, межі сторінок та узгодженість результатів запитів.

\begin{lstlisting}[language={},basicstyle=\ttfamily\small]
scenario inbox pagination

given
  private feed alice<->bob has messages
    1 from alice "p1"
    ...
    11 from alice "p11"

session bob
connect
auth

query inbox peer alice limit 10

expect result items <= 10
expect more

query inbox continue

expect result items
\end{lstlisting}

(Повний набір з DSL-PAGINATION.md: inbox/home/event pagination, continue without initial query, no duplicates, empty page тощо.)

\section{Сценарії присутності}

Сценарії присутності перевіряють стани online/offline, введення повідомлень та інші події присутності.

\begin{lstlisting}[language={},basicstyle=\ttfamily\small]
scenario wildcard offline presence observation

session alice
connect
auth

session bob
connect
auth

session bob
query events peer alice after cursor

session alice
disconnect

session bob
query events peer alice after cursor

expect empty replay
expect event offline
\end{lstlisting}

(Повний набір з DSL-PRESENCE.md: offline/online, typing, multi-session aggregate, federated presence, no rewrite message state тощо.)

\section{Сценарії home}

Сценарії home перевіряють початкове завантаження стану, знімки та їх зв’язок із відтворенням і читанням.

\begin{lstlisting}[language={},basicstyle=\ttfamily\small]
scenario home returns feeds and snapshot after new message

session alice
connect
auth

add bob to roster

session bob
connect
auth

add alice to roster

session alice
send message to bob "m1"

session bob
bootstrap home

expect feeds
expect shared snapshot
expect feeds count <= 10
\end{lstlisting}

(Повний набір з DSL-HOME.md: home bootstrap, shared snapshot, read consistency, policy interaction тощо.)

\section{Сценарії контактів}

Сценарії контактів перевіряють списки контактів і зв’язки між користувачами та їх роль у взаємодіях.

\begin{lstlisting}[language={},basicstyle=\ttfamily\small]
scenario add to roster

session alice
connect
auth

add bob to roster

query roster

expect bob in roster
\end{lstlisting}

(Повний набір з DSL-ROSTER.md: add/remove, one-way/mutual relation, messaging after removal тощо.)

\section{Сценарії згадок}

Сценарії згадок перевіряють формування стану згадок у стрімах подій і їх зв’язок із читанням та видимістю.

\begin{lstlisting}[language={},basicstyle=\ttfamily\small]
scenario mention appears in home after incoming mention message

session alice
connect
auth

add bob to roster

session bob
connect
auth

add alice to roster

session alice
send message to bob {
body: "hi"
mention: bob
}

session bob
bootstrap home

expect feeds
expect mentions
expect shared snapshot
\end{lstlisting}

(Повний набір з DSL-MENTIONS.md: mention in home, read clears mention, replay does not clear, hidden message does not produce mention тощо.)

\section{Сценарії видимості}

Сценарії видимості перевіряють фільтрацію спостережень і обмеження доступу до даних поверх канонічної історії.

\begin{lstlisting}[language={},basicstyle=\ttfamily\small]
scenario message is hidden when clearance is insufficient

given
message m1 has classification topsecret
alice has clearance secret

when alice queries inbox

expect access allowed
expect message m1 hidden
\end{lstlisting}

(Повний набір з DSL-VISIBILITY.md: hidden vs visible, field-level, access denied vs hidden, hidden != deleted тощо.)

\section{Інваріанти та узгодженість}

Інваріанти перевіряють узгодженість між читанням, відтворенням, модерацією, видимістю та іншими властивостями системи.

(Всі invariants з DSL-INVARIANTS.md, ADVANCED.md та cross-layer checks з інших файлів: read cursor unchanged by moderation, ABAC filtering does not change truth, snapshot does not bypass ban, delete overrides edit, no rewrite history тощо.)

\section{Розширення}

Розширення охоплюють спеціалізовані сценарії для автентифікації, політик доступу, пошуку та виявлення можливостей системи.

(Сценарії з DSL-AUTH.md, DSL-ABAC.md, DSL-SEARCH.md, DSL-DISCOVERY.md: auth/resume/renew, ABAC clearance/visibility, search criteria/projection/pagination, discovery capabilities.)

\section{Висновок}

Тестовий набір підсумовується як систематичне покриття поведінки протоколу через сценарії DSL. Він забезпечує повну верифікацію kernel-моделі, invariants та всіх розширень, роблячи протокол Buddha готовим до формальної стандартизації та реалізації.

\bibliographystyle{plain}
\begin{thebibliography}{9}
\bibitem{itu-z120} ITU-T Recommendation Z.120 (02/2011), Message Sequence Chart (MSC).
\bibitem{synrc} synrc/chat repository, ARCH-KERNEL-MAPPINGS.md, 2025--2026.
\bibitem{annexb} ITU-T Z.120 Annex B (04/1998), Formal semantics of Message Sequence Charts.
\end{thebibliography}

\end{document}
